\documentclass[12pt,a4paper]{article}

\usepackage[utf8]{inputenc}
\usepackage[T1]{fontenc}
\usepackage{geometry}
\usepackage{hyperref}
\usepackage{listings}
\usepackage{xcolor}
\usepackage{amsmath}
\usepackage{amssymb}
\usepackage{titlesec}
\usepackage{parskip}
\usepackage{booktabs}
\usepackage{enumitem}
\usepackage{fancyhdr}
\usepackage{graphicx}

\geometry{margin=1in}
\setlength{\headheight}{15pt}

% Code listing style
\definecolor{codebg}{rgb}{0.95,0.95,0.95}
\definecolor{codegreen}{rgb}{0,0.5,0}
\definecolor{codegray}{rgb}{0.4,0.4,0.4}
\definecolor{codepurple}{rgb}{0.5,0,0.5}
\definecolor{codeblue}{rgb}{0,0,0.7}

\lstdefinestyle{cstyle}{
    backgroundcolor=\color{codebg},
    commentstyle=\color{codegreen},
    keywordstyle=\color{codeblue}\bfseries,
    numberstyle=\tiny\color{codegray},
    stringstyle=\color{codepurple},
    basicstyle=\ttfamily\small,
    breakatwhitespace=false,
    breaklines=true,
    captionpos=b,
    keepspaces=true,
    numbers=left,
    numbersep=8pt,
    showspaces=false,
    showstringspaces=false,
    showtabs=false,
    tabsize=4,
    frame=single,
    rulecolor=\color{black!30},
    language=C
}

\lstset{style=cstyle}

% Header / footer
\pagestyle{fancy}
\fancyhf{}
\rhead{ICSP Problem Set}
\lhead{Image Convolution, Rotation \& Cropping}
\cfoot{\thepage}

\hypersetup{
    colorlinks=true,
    linkcolor=blue,
    urlcolor=blue,
    citecolor=blue
}

% ------------------------------------------------------------------
\begin{document}

% Title
\begin{center}
    {\LARGE \textbf{ICSP Problem Set}}\\[0.4em]
    {\large Image Convolution, Rotation, and Cropping}\\[1.2em]
\end{center}

% Drive link — prominently at the top
\begin{center}
    \fbox{%
        \parbox{0.85\linewidth}{%
            \centering
            \textbf{Drive link to code files for each problem}\\[0.4em]
            \url{https://drive.google.com/drive/folders/1oU7uNtzw4zDzvH8hW310cTP0fcf66ybJ?usp=sharing}
        }%
    }
\end{center}

\bigskip
\tableofcontents
\newpage

% ==================================================================
\section{Problem 2-1 — Socialize with Other Pixels \textnormal{[40 points]}}

\subsection*{Background}
Convolution is a fundamental image-processing operation.
A kernel (also called a \emph{filter}) is a small 2-D matrix that is slid over
every pixel of a source image; at each position the kernel values are
multiplied element-wise with the covered pixels and the results are summed to
produce one output pixel.
The original codebase provides a working implementation for a
\(3\times 3\) kernel.

% ------------------------------------------------------------------
\subsection{Part (a) — \texorpdfstring{5\,×\,5}{5x5} Kernel \textnormal{[20 points]}}

\subsubsection*{What was changed}
The loop bounds and the half-width constant were updated from a hard-coded
\texttt{3} to a hard-coded \texttt{5}, and the kernel array was extended to
\(5\times 5\).

Concretely, three things were modified:

\begin{enumerate}[label=\arabic*.]
    \item \textbf{Kernel size constant.}
          \texttt{\#define KERNEL\_SIZE 3} was replaced with
          \texttt{\#define KERNEL\_SIZE 5}.
    \item \textbf{Kernel array.}
          The kernel is now a \(5\times 5\) array initialised (for example)
          with a normalised box-blur:
          \[
              K = \frac{1}{25}\begin{pmatrix}
                  1 & 1 & 1 & 1 & 1 \\
                  1 & 1 & 1 & 1 & 1 \\
                  1 & 1 & 1 & 1 & 1 \\
                  1 & 1 & 1 & 1 & 1 \\
                  1 & 1 & 1 & 1 & 1
              \end{pmatrix}
          \]
    \item \textbf{Half-width.}
          The variable (or expression) that defines how many pixels around the
          centre the kernel reaches changed from \(1\) to \(2\), i.e.\
          \texttt{half = KERNEL\_SIZE / 2}.
\end{enumerate}

\subsubsection*{Code snippet — key changes}

\begin{lstlisting}
/* --- 3x3 original (before) --- */
#define KERNEL_SIZE 3
float kernel[3][3] = { {1/9.0f, 1/9.0f, 1/9.0f},
                        {1/9.0f, 1/9.0f, 1/9.0f},
                        {1/9.0f, 1/9.0f, 1/9.0f} };
int half = 1;  /* KERNEL_SIZE / 2 */

/* --- 5x5 modified (after) --- */
#define KERNEL_SIZE 5
float kernel[5][5] = {
    {1/25.0f, 1/25.0f, 1/25.0f, 1/25.0f, 1/25.0f},
    {1/25.0f, 1/25.0f, 1/25.0f, 1/25.0f, 1/25.0f},
    {1/25.0f, 1/25.0f, 1/25.0f, 1/25.0f, 1/25.0f},
    {1/25.0f, 1/25.0f, 1/25.0f, 1/25.0f, 1/25.0f},
    {1/25.0f, 1/25.0f, 1/25.0f, 1/25.0f, 1/25.0f}
};
int half = 2;  /* KERNEL_SIZE / 2 */
\end{lstlisting}

The convolution loop itself required \emph{no} structural change because it
already used \texttt{half} as its bounds:

\begin{lstlisting}
for (int ky = -half; ky <= half; ky++) {
    for (int kx = -half; kx <= half; kx++) {
        /* accumulate kernel[ky+half][kx+half] * pixel(x+kx, y+ky) */
    }
}
\end{lstlisting}

\subsubsection*{Edge handling}
Pixels within \texttt{half} pixels of the image border are handled by
\emph{clamping}: out-of-bounds accesses are mapped to the nearest valid pixel
coordinate, so the output image retains the same dimensions as the input.

% ------------------------------------------------------------------
\subsection{Part (b) — Arbitrary Odd Kernel Size \textnormal{[20 points]}}

\subsubsection*{What was changed}
The fixed \texttt{\#define} was removed and the kernel size is now accepted as a
run-time parameter (function argument or user input). Two validations are
enforced before any processing begins:

\begin{enumerate}[label=\arabic*.]
    \item \textbf{Odd check.} If the supplied \(n\) is even the program prints an
          error and aborts (or returns an error code).
    \item \textbf{Minimum-size check.} \(n\) must satisfy \(n > 1\); a value of
          \(1\) (identity) or \(0\) is rejected.
\end{enumerate}

\subsubsection*{Code snippet}

\begin{lstlisting}
void convolve(Image *src, Image *dst, float *kernel, int n) {
    /* Validate kernel size */
    if (n % 2 == 0) {
        fprintf(stderr,
            "Error: kernel size must be odd, but got n = %d.\n", n);
        exit(EXIT_FAILURE);
    }
    if (n <= 1) {
        fprintf(stderr,
            "Error: kernel size must be greater than 1, but got n = %d.\n",
            n);
        exit(EXIT_FAILURE);
    }

    int half = n / 2;

    for (int y = 0; y < src->height; y++) {
        for (int x = 0; x < src->width; x++) {
            float acc = 0.0f;
            for (int ky = -half; ky <= half; ky++) {
                for (int kx = -half; kx <= half; kx++) {
                    int sx = clamp(x + kx, 0, src->width  - 1);
                    int sy = clamp(y + ky, 0, src->height - 1);
                    int ki = (ky + half) * n + (kx + half);
                    acc += kernel[ki] * get_pixel(src, sx, sy);
                }
            }
            set_pixel(dst, x, y, (unsigned char)clamp_f(acc, 0, 255));
        }
    }
}
\end{lstlisting}

\subsubsection*{Why only odd sizes?}
An odd kernel size guarantees that the kernel has a well-defined
\emph{centre} pixel at position \(\bigl(\lfloor n/2\rfloor,\,
\lfloor n/2\rfloor\bigr)\).
For an even-sized kernel no such integer centre exists, which would require
sub-pixel alignment and is therefore rejected.

\newpage

% ==================================================================
\section{Problem 2-2 — Take a Left to Go Right \textnormal{[20 points]}}

\subsection*{Background}
Rotating a rectangular image by $90^\circ$ swaps its width and height.
Given an input image of size \(W \times H\) (width \(\times\) height), the
output image has size \(H \times W\).

% ------------------------------------------------------------------
\subsection{Part (a) — Clockwise \texorpdfstring{$90^\circ$}{90 degrees} \textnormal{[10 points]}}

\subsubsection*{Mathematical mapping}
For a clockwise rotation by $90^\circ$ the pixel originally at column \(x\),
row \(y\) (0-indexed, origin at top-left) maps to:
\[
    x' = (H - 1) - y, \qquad y' = x
\]
where \(H\) is the original image height.

\subsubsection*{What was changed}
A new function \texttt{rotate\_cw} was added.
It allocates an output image of dimensions \(H \times W\) and iterates
over every pixel of the input, writing it to the computed destination.

\subsubsection*{Code snippet}

\begin{lstlisting}
Image *rotate_cw(const Image *src) {
    /* Output dimensions are swapped */
    Image *dst = create_image(src->height, src->width);

    for (int y = 0; y < src->height; y++) {
        for (int x = 0; x < src->width; x++) {
            int dx = (src->height - 1) - y;  /* new column */
            int dy = x;                       /* new row    */
            set_pixel(dst, dx, dy, get_pixel(src, x, y));
        }
    }
    return dst;
}
\end{lstlisting}

\subsubsection*{Explanation}
Imagine the image printed on paper and physically rotated clockwise.
The top row of the input becomes the \emph{right-most column} of the output.
Equivalently, the left-most column of the input becomes the \emph{top row} of
the output.

% ------------------------------------------------------------------
\subsection{Part (b) — Counter-clockwise \texorpdfstring{$90^\circ$}{90 degrees} \textnormal{[10 points]}}

\subsubsection*{Mathematical mapping}
For a counter-clockwise rotation by $90^\circ$:
\[
    x' = y, \qquad y' = (W - 1) - x
\]
where \(W\) is the original image width.

\subsubsection*{What was changed}
A complementary function \texttt{rotate\_ccw} was added, mirroring the logic
of \texttt{rotate\_cw} but with the axes transposed in the opposite direction.

\subsubsection*{Code snippet}

\begin{lstlisting}
Image *rotate_ccw(const Image *src) {
    /* Output dimensions are swapped */
    Image *dst = create_image(src->height, src->width);

    for (int y = 0; y < src->height; y++) {
        for (int x = 0; x < src->width; x++) {
            int dx = y;                      /* new column */
            int dy = (src->width - 1) - x;  /* new row    */
            set_pixel(dst, dx, dy, get_pixel(src, x, y));
        }
    }
    return dst;
}
\end{lstlisting}

\subsubsection*{Sanity check}
Applying \texttt{rotate\_cw} followed by \texttt{rotate\_ccw} (or vice versa)
on any image must yield the original image.
This identity was used to verify correctness during development.

\newpage

% ==================================================================
\section{Problem 2-3 — Reshaping Reality \textnormal{[40 points]}}

\subsection*{Background}
This problem extends the codebase with two related operations:
\begin{enumerate}
    \item \textbf{Cropping}: extract a rectangular sub-region from an image.
    \item \textbf{Resizing / stretching}: scale the cropped region to a
          user-specified display size using nearest-neighbour interpolation.
\end{enumerate}

% ------------------------------------------------------------------
\subsection{Part (a) — Cropping \textnormal{[15 points]}}

\subsubsection*{Inputs}
\begin{itemize}
    \item Source image of size \(W \times H\).
    \item Top-left corner \(P = (x_0,\, y_0)\).
    \item Crop width \(w\) and height \(h\).
\end{itemize}

\subsubsection*{What was changed}
A function \texttt{crop} was added.
It validates all parameters (see Section~\ref{sec:safety}), then copies the
rectangular block \([x_0,\, x_0+w) \times [y_0,\, y_0+h)\) from the source
into a freshly allocated output image.

\subsubsection*{Code snippet}

\begin{lstlisting}
Image *crop(const Image *src, int x0, int y0, int w, int h) {
    /* --- safety checks --- */
    if (x0 < 0 || y0 < 0) {
        fprintf(stderr, "Error: crop origin (%d,%d) is out of bounds.\n",
                x0, y0);
        return NULL;
    }
    if (w <= 0 || h <= 0) {
        fprintf(stderr, "Error: crop dimensions (%dx%d) must be positive.\n",
                w, h);
        return NULL;
    }
    if (x0 + w > src->width || y0 + h > src->height) {
        fprintf(stderr,
            "Error: crop region [%d..%d) x [%d..%d) exceeds image "
            "bounds %dx%d.\n",
            x0, x0 + w, y0, y0 + h, src->width, src->height);
        return NULL;
    }

    Image *dst = create_image(w, h);
    for (int y = 0; y < h; y++) {
        for (int x = 0; x < w; x++) {
            set_pixel(dst, x, y, get_pixel(src, x0 + x, y0 + y));
        }
    }
    return dst;
}
\end{lstlisting}

% ------------------------------------------------------------------
\subsection{Part (b) — Stretch / Squeeze \textnormal{[15 points]}}

\subsubsection*{Inputs}
In addition to the cropping parameters, the caller supplies:
\begin{itemize}
    \item \texttt{displayW} — desired output width.
    \item \texttt{displayH} — desired output height.
\end{itemize}

\subsubsection*{What was changed}
A function \texttt{crop\_and\_resize} was added.
It first calls \texttt{crop}, then maps each output pixel
\((x',\, y')\) back to the cropped image using nearest-neighbour
interpolation:
\[
    x_{\text{src}} = \left\lfloor x' \cdot \frac{w}{\texttt{displayW}} \right\rfloor,
    \qquad
    y_{\text{src}} = \left\lfloor y' \cdot \frac{h}{\texttt{displayH}} \right\rfloor
\]

\subsubsection*{Code snippet}

\begin{lstlisting}
Image *crop_and_resize(const Image *src,
                       int x0, int y0, int w, int h,
                       int displayW, int displayH) {
    /* Validate display dimensions */
    if (displayW <= 0 || displayH <= 0) {
        fprintf(stderr,
            "Error: display dimensions (%dx%d) must be positive.\n",
            displayW, displayH);
        return NULL;
    }

    Image *cropped = crop(src, x0, y0, w, h);
    if (!cropped) return NULL;  /* propagate crop error */

    Image *dst = create_image(displayW, displayH);

    for (int dy = 0; dy < displayH; dy++) {
        for (int dx = 0; dx < displayW; dx++) {
            /* Nearest-neighbour mapping back into the cropped region */
            int sx = (int)((double)dx * w / displayW);
            int sy = (int)((double)dy * h / displayH);
            /* Clamp to [0, w-1] x [0, h-1] for safety */
            sx = sx < w ? sx : w - 1;
            sy = sy < h ? sy : h - 1;
            set_pixel(dst, dx, dy, get_pixel(cropped, sx, sy));
        }
    }

    free_image(cropped);
    return dst;
}
\end{lstlisting}

% ------------------------------------------------------------------
\subsection{Part (c) — Safety Checks and Edge Cases}
\label{sec:safety}

The following table summarises every guard implemented across
\texttt{crop} and \texttt{crop\_and\_resize}.

\begin{center}
\begin{tabular}{@{}lll@{}}
    \toprule
    \textbf{Check} & \textbf{Condition tested} & \textbf{Action} \\
    \midrule
    Negative origin
        & \(x_0 < 0\) or \(y_0 < 0\)
        & Print error, return \texttt{NULL} \\
    Non-positive crop size
        & \(w \le 0\) or \(h \le 0\)
        & Print error, return \texttt{NULL} \\
    Crop exceeds image
        & \(x_0 + w > W\) or \(y_0 + h > H\)
        & Print error, return \texttt{NULL} \\
    Non-positive display size
        & \(\texttt{displayW} \le 0\) or \(\texttt{displayH} \le 0\)
        & Print error, return \texttt{NULL} \\
    Nearest-neighbour clamp
        & Computed \(sx \ge w\) or \(sy \ge h\)
        & Clamp to \(w-1\) / \(h-1\) \\
    Null pointer propagation
        & \texttt{crop} returns \texttt{NULL}
        & \texttt{crop\_and\_resize} returns \texttt{NULL} \\
    \bottomrule
\end{tabular}
\end{center}

Additional edge cases noted:

\begin{itemize}
    \item \textbf{Point crop} (\(w = 1,\, h = 1\)): produces a
          \(1 \times 1\) image — valid and handled correctly.
    \item \textbf{Full-image crop} (\(x_0 = 0,\, y_0 = 0,\, w = W,\, h = H\)):
          produces an exact copy — verified against the original.
    \item \textbf{Identity resize} (\(\texttt{displayW} = w,\,
          \texttt{displayH} = h\)): output equals the cropped image
          pixel-for-pixel.
    \item \textbf{Extreme up-scaling}: e.g.\ a \(1 \times 1\) crop stretched
          to \(1920 \times 1080\) — every output pixel maps to the single source
          pixel; clamping ensures no out-of-bounds access.
    \item \textbf{Extreme down-scaling}: e.g.\ a large crop compressed to
          \(1 \times 1\) — only the top-left pixel of the crop is retained.
\end{itemize}

% ==================================================================
\section*{Summary of All Modifications}

\begin{center}
\begin{tabular}{@{}llp{7cm}@{}}
    \toprule
    \textbf{Problem} & \textbf{Function / change} & \textbf{Key modification} \\
    \midrule
    2-1(a) & \texttt{convolve} (5\,×\,5)
            & Extended kernel array and \texttt{half} from 1 to 2. \\
    2-1(b) & \texttt{convolve} (n\,×\,n)
            & Kernel size becomes a parameter; even/small values rejected. \\
    2-2(a) & \texttt{rotate\_cw}
            & Pixel \((x,y)\) maps to \((H-1-y,\;x)\). \\
    2-2(b) & \texttt{rotate\_ccw}
            & Pixel \((x,y)\) maps to \((y,\;W-1-x)\). \\
    2-3(a) & \texttt{crop}
            & Copies rectangle \([x_0, x_0+w) \times [y_0, y_0+h)\). \\
    2-3(b) & \texttt{crop\_and\_resize}
            & Nearest-neighbour scale of crop to \(\texttt{displayW}\times\texttt{displayH}\). \\
    \bottomrule
\end{tabular}
\end{center}

\end{document}
